\section{Chapter 5 supplementary exercises}
These exercises connect eigenvectors, similarity, polynomial expressions
in matrices, and diagonalization. Keep three distinctions in view:
row equivalence does not preserve eigenvalues; similarity does.
An eigenvalue may be zero, but an eigenvector must be nonzero.
A repeated eigenvalue does not by itself decide diagonalizability:
the total number of independent eigenvectors decides it.
All matrices are real unless complex roots or complex diagonalization
are explicitly being considered.

\exercise{5.S.1}{Twenty-four statements with proofs or counterexamples}
\begin{enumerate}[label=(\alph*)]
\item \textbf{True.} If $Ax=x$ and $x\ne0$, multiplication by
$A^{-1}$ gives $A^{-1}x=x$.
\item \textbf{False.} $\begin{bmatrix}1&1\\0&1\end{bmatrix}$ is
invertible and row equivalent to $I$, but has only one independent
eigenvector, since its $1$-eigenspace is $\Span\{(1,0)^T\}$.
\item \textbf{True.} A zero column makes $Ae_j=0$ for some nonzero
coordinate vector. A zero row makes the matrix singular as well;
by the Invertible Matrix Theorem its kernel contains a nonzero vector.
\item \textbf{False.} The scalar matrix $[2]$ has eigenvalue $2$,
whereas its square $[4]$ has only eigenvalue $4$.
\item \textbf{True.} $Ax=\lambda x$ implies
$A^2x=\lambda^2x$, with the same nonzero $x$.
\item \textbf{True.} An invertible $A$ has $\lambda\ne0$;
$Ax=\lambda x$ gives $A^{-1}x=\lambda^{-1}x$.
\item \textbf{False.} The zero matrix has eigenvalue zero.
\item \textbf{True.} Nonzeroness is part of the eigenvector definition;
otherwise $A0=\lambda0$ would make every scalar an eigenvalue.
\item \textbf{False.} For $I_2$, the independent vectors $e_1,e_2$
have the same eigenvalue $1$.
\item \textbf{True.} Similarity gives
$\det(P^{-1}(A-\lambda I)P)=\det(A-\lambda I)$, so the
characteristic polynomials, including multiplicities, agree.
\item \textbf{False.} For $A=\operatorname{diag}(1,2)$ and
$B=\begin{bmatrix}1&1\\0&2\end{bmatrix}$ the eigenvalues agree
and the matrices are similar (both have independent eigenbases),
but $e_2$ is an eigenvector of $A$ and not of $B$.
\item \textbf{False.} For $A=\operatorname{diag}(1,2)$,
$e_1+e_2$ is not an eigenvector. Even within one eigenspace,
the sum of $v$ and $-v$ is zero and hence not an eigenvector.
\item \textbf{False.} All diagonal entries of an upper triangular
matrix are eigenvalues, including zeros.
\item \textbf{True.} $\det(A^T-\lambda I)=
\det((A-\lambda I)^T)=\det(A-\lambda I)$.
\item \textbf{False.} $I_5$ has only one distinct eigenvalue but is
already diagonal; its five coordinate vectors form an eigenbasis.
\item \textbf{True.} The real rotation matrix
$\begin{bmatrix}0&-1\\1&0\end{bmatrix}$ has eigenvalues $\pm i$
and no real eigenvectors.
\item \textbf{False.} $\operatorname{diag}(0,1)$ is diagonalizable
but its first column is zero, so its columns are dependent.
\item \textbf{True.} If $Av=\lambda v=\mu v$ and $v\ne0$,
then $(\lambda-\mu)v=0$ forces $\lambda=\mu$.
\item \textbf{False.} Invertibility and diagonalizability are different.
The non-diagonalizable invertible example in (b) refutes one direction;
$\operatorname{diag}(0,1)$ refutes the reverse direction.
\item \textbf{True.} If $Ae_j=\lambda_je_j$, column $j$ has
only its diagonal entry possibly nonzero, so $A$ is diagonal.
\item \textbf{True.} If $A=SBS^{-1}$ and $B=PDP^{-1}$,
then $A=(SP)D(SP)^{-1}$, a diagonalization.
\item \textbf{True.} $A^{-1}(AB)A=BA$, so the products are
similar. In fact invertibility of $A$ alone suffices.
\item \textbf{False.} The zero matrix has $n$ independent
eigenvectors but is not invertible.
\item \textbf{True.} Diagonalizability supplies an eigenvector basis,
and the definition of a basis says every vector is a linear
combination of its vectors (over $\R$ for real diagonalization here).
\end{enumerate}

\exercise{5.S.2}{Move an eigenvector through the right-hand factor}
Suppose $(AB)x=\lambda x$ and $Bx\ne0$. Associativity gives
\[(BA)(Bx)=B(ABx)=B(\lambda x)=\lambda(Bx).\]
Because $Bx$ is nonzero by assumption, it is an eigenvector of
$BA$ for the same eigenvalue. The condition $Bx\ne0$ is what
prevents the zero vector from being incorrectly called an eigenvector.

\exercise{5.S.3}{Apply two polynomials to an eigenvector}
\begin{enumerate}[label=(\alph*)]
\item $(5I-A)x=5x-Ax=(5-\lambda)x$. Thus $x$ remains a
nonzero eigenvector, with eigenvalue $5-\lambda$.
\item First $A^2x=A(\lambda x)=\lambda^2x$. Consequently
$(5I-3A+A^2)x=(5-3\lambda+\lambda^2)x$.
The new eigenvalue is $5-3\lambda+\lambda^2$.
\end{enumerate}

\exercise{5.S.4}{Induct on the matrix power}
For $m=1$ the assertion is the assumed eigenvector equation.
Suppose $A^mx=\lambda^mx$. Multiplying by $A$ gives
$A^{m+1}x=A(\lambda^mx)=\lambda^mAx=\lambda^{m+1}x$.
The same $x$ remains nonzero at every stage, so each displayed
scalar is genuinely an eigenvalue of the corresponding power.
This proves the claim for every positive integer $m$.

\exercise{5.S.5}{Polynomial evaluation preserves eigenvectors}
Exercise 5.S.4 gives $A^jx=\lambda^jx$ for every positive power;
for $j=0$ the same equality is $Ix=x$. Therefore
\[p(A)x=\sum_{j=0}^n c_jA^jx
=\sum_{j=0}^n c_j\lambda^jx=p(\lambda)x.\]
Since $x\ne0$, $p(\lambda)$ is an eigenvalue of $p(A)$.
Different eigenvalues of $A$ can map to the same polynomial value;
that does not invalidate the conclusion.

\exercise{5.S.6}{Diagonalize polynomial expressions using the original basis}
\begin{enumerate}[label=(\alph*)]
\item With $A=PDP^{-1}$, neighboring factors $P^{-1}P$ cancel
in $A^2$. Since $I=PIP^{-1}$,
$5I-3A+A^2=P(5I-3D+D^2)P^{-1}$.
At diagonal entries $2,7$, the polynomial has values $3,33$.
Thus $B=P\operatorname{diag}(3,33)P^{-1}$ is a diagonalization.
\item The same cancellation gives $A^j=PD^jP^{-1}$ for every
power. Adding the finitely many terms gives
$p(A)=Pp(D)P^{-1}$, with diagonal entries $p(2),p(7)$.
This remains a diagonalization even if those two values coincide.
\end{enumerate}

\exercise{5.S.7}{Cayley--Hamilton for a diagonalizable matrix}
Write $A=PDP^{-1}$. Every diagonal entry $\lambda_j$ of $D$
is an eigenvalue of $A$ and hence a root of its characteristic
polynomial $p$. By the polynomial identity in Exercise 5.S.6,
\[p(A)=Pp(D)P^{-1}=P\operatorname{diag}(p(\lambda_1),\ldots,
p(\lambda_n))P^{-1}=P0P^{-1}=0.\]
This proves the theorem under the stated diagonalizability
assumption; it does not pretend that every square matrix is
already known to be diagonalizable.

\exercise{5.S.8}{One eigenvalue plus diagonalizability forces a scalar matrix}
\begin{enumerate}[label=(\alph*)]
\item If the only eigenvalue is $\lambda$ with multiplicity $n$,
a diagonalization has $D=\lambda I$. Therefore
$A=P(\lambda I)P^{-1}=\lambda PP^{-1}=\lambda I$.
\item The given $\begin{bmatrix}3&1\\0&3\end{bmatrix}$ has only
eigenvalue $3$, twice, but is not $3I$. By (a), it cannot be
diagonalizable. Directly, $(A-3I)x=0$ forces $x_2=0$,
leaving only one independent eigenvector.
\end{enumerate}

\exercise{5.S.9}{An eigenvalue inside the unit disk cannot equal one}
If $I-A$ were singular, a nonzero $v$ would satisfy $(I-A)v=0$.
Then $Av=v$, so $1$ would be an eigenvalue of $A$.
Its modulus is $1$, contradicting the strict bound on every
eigenvalue. Hence $I-A$ is invertible. This argument requires no
diagonalizability assumption.

\exercise{5.S.10}{Every eigenmode decays and therefore the whole matrix does}
Use $A=PDP^{-1}$. Its powers satisfy $A^k=PD^kP^{-1}$.
Every diagonal entry of $D^k$ is $\lambda_j^k\to0$ because
$|\lambda_j|<1$. Thus each entry of $PD^kP^{-1}$ is a finite
sum of fixed constants times quantities tending to zero, so it also
tends to zero. Equivalently, apply the mode decomposition to each
coordinate vector $e_j$; then $A^ke_j\to0$, which says each
column of $A^k$ tends to zero.
\answer{$A^k\to0$ entry by entry.}

\exercise{5.S.11}{Invariant lines and real eigenvectors}
\begin{enumerate}[label=(\alph*)]
\item Every $x\in H$ has the form $x=cu$. Then
$Ax=cAu=c\lambda u\in H$. This is exactly invariance.
\item Choose a nonzero $u$ spanning the one-dimensional invariant
subspace $K$. Invariance gives $Au\in K$, so $Au=\lambda u$
for some scalar $\lambda$. Since $u\ne0$, it is an eigenvector.
\end{enumerate}

\exercise{5.S.12}{Factor a block triangular characteristic polynomial}
In the permutation expansion of
$G=\begin{bmatrix}A&X\\0&B\end{bmatrix}$, a nonzero term must
assign the bottom block rows to bottom block columns; otherwise it
uses a zero from the lower-left block. The remaining top rows use
top block columns. The products and signs therefore separate into
the determinant sums for $A$ and $B$, giving $\det G=\det A\det B$.
The same argument applied to
$G-\lambda I=\begin{bmatrix}A-\lambda I&X\\0&B-\lambda I\end{bmatrix}$
gives the product of the two characteristic polynomials. No inverse
of either diagonal block is required.

\exercise{5.S.13}{Read one eigenvalue and solve the remaining block}
The upper-left block is $[3]$ and the lower-right block is
$B=\begin{bmatrix}5&-2\\-4&3\end{bmatrix}$.
Its characteristic polynomial is
$(\lambda-5)(\lambda-3)-8=\lambda^2-8\lambda+7
=(\lambda-1)(\lambda-7)$.
By Exercise 5.S.12 the complete eigenvalue list is $3,1,7$,
each with multiplicity one. The entries in the upper-right block
have no effect on this list.

\exercise{5.S.14}{Combine the spectra of two diagonal blocks}
The blocks are $B=\begin{bmatrix}1&5\\2&4\end{bmatrix}$ and
$C=\begin{bmatrix}-7&-4\\3&1\end{bmatrix}$.
Their characteristic polynomials are
\[\det(\lambda I-B)=\lambda^2-5\lambda-6=(\lambda-6)(\lambda+1),\]
\[\det(\lambda I-C)=\lambda^2+6\lambda+5=(\lambda+1)(\lambda+5).\]
Multiply them by the block result. The eigenvalues are $6,-5$,
each once, and $-1$ twice. Algebraic multiplicity is counted even
when an eigenvalue occurs in both blocks.

\exercise{5.S.15}{Split into the constant vector and coordinate-sum-zero space}
Let $\mathbf1=(1,\ldots,1)^T$. Then $J\mathbf1=n\mathbf1$,
so $A\mathbf1=[a+(n-1)b]\mathbf1$.
For $z$ with coordinate sum zero, $Jz=0$, and hence
$Az=(a-b)z$. The coordinate-sum-zero subspace has dimension
$n-1$; together with $\mathbf1$ it gives a basis of $\R^n$.
Thus the characteristic polynomial is
\[(\lambda-(a-b))^{n-1}(\lambda-(a+(n-1)b)).\]
When $b\ne0$ and $n\ge2$, the two eigenvalues are distinct, with
multiplicities $n-1$ and $1$. If $b=0$, they coincide at $a$ and
the combined multiplicity is $n$; for $n=1$ only the value $a$
is present. These qualifications prevent double-counting a coincident
value.

\exercise{5.S.16}{Apply the common-diagonal formula twice}
For the $3\times3$ matrix, $a=1,b=2$. Thus $a-b=-1$ has
multiplicity $2$, and $a+2b=5$ has multiplicity $1$.
For the $5\times5$ matrix, $a=7,b=3$. Thus $a-b=4$ has
multiplicity $4$, and $a+4b=19$ has multiplicity $1$.
These multiplicities add to the dimension in each case.

\exercise{5.S.17}{The trace and determinant decide reality in dimension two}
For a real matrix $A=[a_{ij}]$, expand
\[\det(\lambda I-A)
=(\lambda-a_{11})(\lambda-a_{22})-a_{12}a_{21}
=\lambda^2-(\operatorname{tr}A)\lambda+\det A.\]
A real quadratic has both roots real exactly when its discriminant
is nonnegative. Here that condition is
$(\operatorname{tr}A)^2-4\det A\ge0$, equivalently
\[\det A\le\left(\frac{\operatorname{tr}A}{2}\right)^2.\]
Equality includes a repeated real eigenvalue; it does not by itself
assert diagonalizability.

\exercise{5.S.18}{Compute the limiting spectral projection}
The matrix has eigenvalues $1$ and $3/5$, with eigenvectors
$(-1,2)^T$ and $(-3,2)^T$. Distinct eigenvalues give a basis.
The projection onto the $1$-eigenspace along the other eigenspace is
\[E=\frac{A-(3/5)I}{1-3/5}
=\begin{bmatrix}-1/2&-3/4\\1&3/2\end{bmatrix}.\]
Indeed $E$ acts as identity on the first eigenvector and as zero
on the second, so on the whole eigenbasis
\[A^k=E+(3/5)^k(I-E).\]
The second term tends to zero. Therefore $A^k$ tends to the
stated matrix $E$, whose negative entries cause no contradiction:
the original $A$ already has a negative off-diagonal entry.

\exercise{5.S.19}{A quadratic companion matrix}
For $p(t)=6-5t+t^2$, the companion matrix in the book's convention is
\[C_p=\begin{bmatrix}0&1\\-6&5\end{bmatrix}.\]
Its characteristic polynomial is
$\det(\lambda I-C_p)=\det\begin{bmatrix}\lambda&-1\\6&\lambda-5\end{bmatrix}
=\lambda^2-5\lambda+6=p(\lambda)$.
If using $\det(C_p-\lambda I)$ instead, the sign is the same in
even dimension.

\exercise{5.S.20}{A cubic companion matrix and its determinant}
The coefficients $(-24,26,-9,1)$ give
\[C_p=\begin{bmatrix}0&1&0\\0&0&1\\24&-26&9\end{bmatrix}.\]
Expand $\det(\lambda I-C_p)$ along the first row:
\[\lambda\det\begin{bmatrix}\lambda&-1\\26&\lambda-9\end{bmatrix}
-24=\lambda^3-9\lambda^2+26\lambda-24.
\]
This equals $(\lambda-2)(\lambda-3)(\lambda-4)=p(\lambda)$.
For the book's $\det(C_p-\lambda I)$ convention the polynomial
is its negative, as the dimension is odd.

\exercise{5.S.21}{Inductively compute the companion characteristic polynomial}
For $n=2$, direct expansion gives
$\det\begin{bmatrix}-\lambda&1\\-a_0&-a_1-\lambda\end{bmatrix}
=\lambda^2+a_1\lambda+a_0=(-1)^2p(\lambda)$.
Suppose the identity holds in dimension $n-1$.
Expand $C_p-\lambda I$ down its first column. Only the top and
bottom entries can contribute. The top contribution is $-\lambda$
times the determinant of the companion-form block associated with
$q(t)=a_1+a_2t+\cdots+a_{n-1}t^{n-2}+t^{n-1}$.
By induction this determinant is $(-1)^{n-1}q(\lambda)$.
The bottom contribution is $(-a_0)(-1)^{n+1}$ times a determinant
of $1$; the latter minor is triangular with diagonal ones.
Thus the total is
\[(-\lambda)(-1)^{n-1}q(\lambda)+(-1)^na_0
=(-1)^n[\lambda q(\lambda)+a_0]=(-1)^np(\lambda).\]
This completes the induction for all $n\ge2$.

\exercise{5.S.22}{A polynomial root supplies an explicit companion eigenvector}
\begin{enumerate}[label=(\alph*)]
\item The companion matrix is
$C_p=\begin{bmatrix}0&1&0\\0&0&1\\-a_0&-a_1&-a_2\end{bmatrix}$.
\item Since $p(\lambda)=0$,
$\lambda^3=-a_0-a_1\lambda-a_2\lambda^2$. Therefore
\[C_p(1,\lambda,\lambda^2)^T
=(\lambda,\lambda^2,-a_0-a_1\lambda-a_2\lambda^2)^T
=(\lambda,\lambda^2,\lambda^3)^T
=\lambda(1,\lambda,\lambda^2)^T.\]
The first coordinate is $1$, so the vector is nonzero even if
$\lambda=0$; it is genuinely an eigenvector.
\end{enumerate}

\exercise{5.S.23}{A Vandermonde matrix is an eigenvector matrix}
Each column of $V$ is the eigenvector from Exercise 5.S.22 for
one root $\lambda_j$. Distinct roots give distinct eigenvalues,
so the corresponding eigenvectors are independent. Thus $V$ is
invertible without needing to compute its inverse.
The three eigenvector equations combine into
$C_pV=V\operatorname{diag}(\lambda_1,\lambda_2,\lambda_3)$.
Multiply on the left by $V^{-1}$ to obtain the requested diagonal
matrix. Complex roots are allowed if matrices and independence are
interpreted over $\C$.

\exercise{5.S.24}{The companion-matrix idea behind a polynomial-root command}
The basic method documented for MATLAB's \texttt{roots} is to
construct a companion matrix and compute its eigenvalues. After
removing leading zero coefficients and normalizing the leading
nonzero coefficient, the polynomial is monic. Exercise 5.S.21
shows that the companion matrix's eigenvalues are exactly its roots.
Numerical eigenvalue algorithms, typically using QR-based methods,
then approximate those values. The companion matrix used by a
program may be transposed or written in a different convention;
its characteristic polynomial still supplies the same roots.
The MATLAB reference is
\url{https://www.mathworks.com/help/matlab/ref/roots.html}.
This explains the algorithmic idea, not an assertion that every
computed root has exact arithmetic accuracy: companion matrices can
be sensitive to coefficient rounding, especially near repeated roots.

\exercise{5.S.25}{A repeated eigenvalue prevents an exact diagonalization}
For the supplied matrix, exact symbolic computation gives
\[\det(\lambda I-A)=(\lambda-1)(\lambda-2)^2.\]
Its eigenspaces are
\[E_1=\Span\{(1,-2,2)^T\},\qquad
E_2=\Span\{(2,-5,3)^T\}.\]
Each has dimension one, so there are only two independent eigenvectors
for a three-dimensional matrix. Exact diagonalization is impossible.
To see what happens to the requested checks, choose
\[D=\operatorname{diag}(1,2,2),\quad
P=\begin{bmatrix}1&2&2\\-2&-5&-5\\2&3&3\end{bmatrix}.\]
Then $AP-PD=0$, but $\det P=0$ because columns 2 and 3 coincide.
Thus $PDP^{-1}$ is undefined. A zero eigenvector-equation residual
alone does not make $P$ an invertible change of basis.
Floating-point eigenvalue commands may split the repeated root
slightly and return nearly parallel eigenvectors. An attempted inverse
then amplifies rounding errors; exact ranks resolve the issue.
The following executed symbolic checks reproduce the conclusion:
\begin{verbatim}
from sympy import Matrix, eye, symbols, factor
A = Matrix([[-3,-2,0],[14,7,-1],[-6,-3,1]])
z = symbols('z')
print(factor(A.charpoly(z).as_expr()))
print((A-eye(3)).nullspace())
print((A-2*eye(3)).nullspace())
P = Matrix([[1,2,2],[-2,-5,-5],[2,3,3]])
D = Matrix.diag(1,2,2)
assert A*P-P*D == Matrix.zeros(3)
assert P.det() == 0
\end{verbatim}

\exercise{5.S.26}{A nonzero nilpotent matrix cannot be diagonalized}
For the given $4\times4$ matrix, exact calculation gives
\[\det(\lambda I-A)=\lambda^4,\qquad
\Nul A=\Span\{(1,2,1,0)^T\}.\]
The only eigenvalue is zero, with algebraic multiplicity four but
only one independent eigenvector. Thus the matrix is not diagonalizable.
In fact $A^4=0$ and $A^3\ne0$, another exact check of its nonzero
nilpotent structure.
An eigenvector matrix for $D=0_4$ can only have columns proportional
to $(1,2,1,0)^T$. Such a $P$ has rank at most one: $AP-PD=0$
is possible, but $P^{-1}$ is impossible. Were an invertible $P$
to exist, $PDP^{-1}=0$ would contradict the given nonzero $A$.
A floating-point eigenvalue routine can again produce small nonzero
numbers around the repeated zero root and unreliable eigenvectors.
The exact calculation that settles the issue is:
\begin{verbatim}
A = Matrix([[-8,5,-2,0],[-5,2,1,-2],
            [10,-8,6,-3],[3,-2,1,0]])
assert A.charpoly().as_expr().expand() == A.charpoly().gen**4
assert len(A.nullspace()) == 1
assert A**4 == Matrix.zeros(4)
assert A**3 != Matrix.zeros(4)
\end{verbatim}
No external data are needed, and the failure is a proved mathematical
obstruction, not a software or setup failure.
